\documentclass[11pt]{article}
\usepackage[a4paper,margin=23mm]{geometry}
\usepackage[T1]{fontenc}
\usepackage{lmodern,amsmath,amssymb,amsthm,mathtools,microtype,booktabs}
\usepackage[colorlinks=true,linkcolor=blue,citecolor=blue,urlcolor=blue]{hyperref}
\newcommand{\EE}{\mathbb E}
\newcommand{\tr}{\operatorname{tr}}
\newcommand{\Om}{\Omega_O}
\newcommand{\nn}{\hat n}
\newcommand{\KK}{\mathcal K}
\newcommand{\ff}[2]{(#1)_{\underline{#2}}}
\newtheorem{theorem}{Theorem}
\newtheorem{proposition}{Proposition}
\title{A Slater-adapted oscillator in complex double-scaled SYK\\
\large Same-disorder polarization, all fixed vacuum words, and the first Hamiltonian defect}
\author{Alan Day \and Chengheng\thanks{AI research assistant (OpenAI).}}
\date{7 November 2026\quad Local research draft}
\begin{document}\maketitle
\begin{abstract}
We select the pure particle--hole creation block of the original complex
SYK Hamiltonian relative to a half-filled Slater state. All fixed ordered
vacuum moments of its normalized creation and annihilation operators
converge to a single dressed bosonic oscillator,
$M=q^{\hat n/2}b$, $[b,b^\dagger]=1$, with $q=e^{-4\lambda}$.
The proof counts all forward Gaussian pairings and bounds the finite-size
backward pairings rather than discarding them. The construction implements
the state-adapted strategy of recent pure-state DSSYK work in a
charge-conserving, same-disorder model. Microscopic occupation-number
undressing gives asymptotic canonical commutators on a fixed polynomial
core. The original Hamiltonian has
additional directions: its first action has a Poisson$(2\lambda)$
particle--hole deficit and only a fraction $\sqrt q$ of its squared norm
lies in the pure-creation direction. This identifies a concrete next
calculation, the joint $(H,M,M^\dagger)$ law.
\end{abstract}

\section{Literature-led question and microscopic conventions}
The Susskind static-patch program motivates deriving particle operators
from the hologram. The quantitative singlet/meson correspondence remains
a target in \cite[Sec.~4.3]{MSS}. Here we pursue the state and polarization
input following the state-adapted construction of Rajgadia--Xu
\cite[Sec.~3.2]{RX}. Their dressed oscillator representation is a
predecessor, not a new general construction claimed in this note. Our
task is to derive a microscopic realization with the original complex
SYK couplings and then test its Hamiltonian closure.

Use $N$ complex fermions, $\{c_i,c_j^\dagger\}=\delta_{ij}$, and
\begin{equation}\label{ham}
 H_N=\sqrt{w_N}\sum_{|A|=|B|=p}J_{AB}C_A^\dagger C_B,
 \quad w_N=\frac{4^pJ_s^2}{\lambda_N}\left(\frac{p!}{N^p}\right)^2,
 \quad \lambda_N=\frac{p^2}{N}.
\end{equation}
$C_A^\dagger$ is the product of creators in increasing site order and
$C_A=(C_A^\dagger)^\dagger$. The centered Hermitian Gaussian matrix has
$\EE J_{AB}J_{CD}=\delta_{AD}\delta_{BC}$; overlapping $A,B$ are included.
All operators below use this same $J$. Take even $N$, $m=N/2$,
$1\le p\le m$, $p\to\infty$, $\lambda_N\to\lambda>0$, fixed $J_s>0$,
and $\hbar=1$.
The Hamiltonian uses string-energy units. All word lengths are fixed
before this limit; no long-time or growing-excitation limit is taken.

\newpage
\section{Keep the reference polarization in the operator}
Fix an occupied set $O$, $|O|=m$, and its empty complement $E$,
$|E|=a=N-m$. Write $\Om=|O\rangle$ and
\begin{equation}\label{state}
 \phi_{O,N}(X)=\EE_J\langle O|X(J)|O\rangle.
\end{equation}
This is an annealed vector state; it is not the full-Fock trace of the
preceding bilocal calculations. The reference state is not asserted
to be stationary under $H_N$.

\paragraph{Why changing only the state is insufficient.}
For a jointly $U(N)$-covariant word $W(J)$, $\EE W$ commutes with
the $U(N)$ action on $\wedge^m\mathbb C^N$. Irreducibility and Schur's
lemma give the exact finite-size identity
\begin{equation}\label{schur}
 \phi_{O,N}(W)=\EE_J\tr_m W,
 \qquad \tr_m=\binom Nm^{-1}\operatorname{Tr}_{\mathcal H_m}.
\end{equation}
No Haar average of $|O\rangle$ is needed. In particular all annealed
$H$-only moments are unchanged. This explains why the operator must
retain information about $O$.

Define the original Hamiltonian's pure-excitation block and its normalization:
\begin{align}
 V_N^\dagger&=\sqrt{w_N}
 \sum_{\substack{A\subset E,\ B\subset O\\|A|=|B|=p}}
 J_{AB}C_A^\dagger C_B,\label{block}\\
 L_N&=\binom ap\binom mp,\qquad
 \sigma_N^2=w_NL_N,\qquad
 M_N^\dagger=V_N^\dagger/\sigma_N,\quad M_N=(M_N^\dagger)^\dagger.
 \label{normalize}
\end{align}
The selected $J_{AB}$ are independent circular complex Gaussians of
unit variance. Their conjugates occur in $V_N$. These are a subfamily
of the original disorder, not an independent probe ensemble.
Every selected monomial preserves the physical total charge $Q$,
while creating $p$ particles in $E$ and $p$ holes in $O$. Exactly,
\begin{equation}\label{vac}
 M_N\Om=0,\qquad
 \phi_{O,N}(M_NM_N^\dagger)=1,\qquad
 \phi_{O,N}(M_N^\dagger M_N)=0.
\end{equation}
The occupation projector is additional reference data. Equation
\eqref{schur} does not apply to a word that keeps this projector fixed
while rotating the fermions and couplings.

This choice differs from taking the empty physical complex-fermion
vacuum: the number-conserving $H_N$ annihilates that vacuum. In the
Majorana KM construction of \cite{RX}, the analogous spin-flip number
is not the conserved charge of \eqref{ham}. The two microscopic
constructions therefore require separate pairing derivations.

\newpage
\section{All fixed words: active intervals and backward corrections}
Read a word in $M_N,M_N^\dagger$ from right to left. A creation raises
the particle--hole block count by one; an annihilation lowers it by one.
A balanced word gives a Dyck path if every intermediate height is
nonnegative. For a creation at height $h$, set
\begin{equation}\label{ratio}
 R_N(h)=\frac{\binom{a-hp}{p}\binom{m-hp}{p}}{L_N},
 \qquad q=e^{-4\lambda},
\end{equation}
with out-of-range binomials zero.

\begin{theorem}[Fixed-word limit]\label{wordthm}
For each fixed word $w$, $\phi_{O,N}(w(M_N,M_N^\dagger))$ tends to
the vacuum moment of the weighted shift
\begin{equation}\label{shift}
 M|n\rangle=\sqrt{nq^{n-1}}\,|n-1\rangle,\qquad
 M^\dagger|n\rangle=\sqrt{(n+1)q^n}\,|n+1\rangle
\end{equation}
on $\ell^2(\mathbb N_0)$. For a balanced Dyck word of length $2k$,
its finite-size forward-pairing contribution is exactly
\begin{equation}\label{forward}
 F_N(w)=\left(\prod_{\text{annihilations}}h\right)
                  \prod_{\text{creations}}R_N(h).
\end{equation}
The remaining contribution satisfies
\begin{equation}\label{backbound}
 |\phi_{O,N}(w)-F_N(w)|\le k!\rho_N(k),\quad
 \rho_N(k)=\frac{\binom{kp}{p}}{\binom ap}
           \frac{\binom{kp}{p}}{\binom mp}\longrightarrow0.
\end{equation}
An unbalanced or non-Dyck word has zero vacuum moment already at
finite size.
\end{theorem}

\paragraph{Forward pairings.}
Complex Gaussian Wick expansion matches each creation coefficient
to an annihilation coefficient. A pair is forward if its creation acts
first. For an all-forward matching, a label is active between its two
endpoints. At each creation, both its particle set and its hole set
must avoid every active set. This condition is necessary and sufficient
for a nonzero CAR word. Legal active sets are pairwise disjoint; each
composite monomial has even degree $2p$, so active monomials commute.
Each annihilation removes its own active monomial with positive
amplitude. Thus every allowed assignment contributes exactly $+1$.

With $h$ labels active, the next creation has exactly
$\binom{a-hp}{p}\binom{m-hp}{p}$ choices. A closing endpoint has $h$
choices of partner. Normalizing every pair by $L_N$ proves
\eqref{forward}. Both nested and crossed intervals impose exclusion
when simultaneously active. Counting only crossing chords would give
the wrong law.

\newpage
\subsection*{Backward pairings are small, not identically zero}
In a matching with a backward pair, choose the first annihilation
whose matched creation has not yet acted. Before this endpoint, every
nonzero action is an all-forward history. Its particle and hole sets
must each lie inside the unions already excited, each of size at most
$kp$. This newly encountered Gaussian label is independent of all
previous labels. Its conditional allowed fraction is at most
$\rho_N(k)$. All later CAR amplitudes have absolute value at most one.
There are at most $k!$ matchings, proving \eqref{backbound}.
For fixed $k$ and sufficiently large $N$,
\begin{equation}
 \rho_N(k)\le
 \left(\frac{kp}{a-p+1}\frac{kp}{m-p+1}\right)^p\longrightarrow0.
\end{equation}
This argument also allows different Gaussian labels to take the same
site subsets; they are not excluded by hand.

For example, the alternating fourth moment has the finite correction
\begin{equation}\label{backexample}
 \phi_{O,N}(M_NM_N^\dagger M_NM_N^\dagger)=1+\frac1{L_N}.
\end{equation}
The extra matching repeats the same particle and hole sets. It is an
explicit check against replacing all finite-size pairings by forward
ones prematurely.

\paragraph{Taking the limit.}
For each fixed $h$,
\begin{equation}
 \log R_N(h)
 =-hp^2\left(\frac1a+\frac1m\right)+O_h(p^3/N^2)
 \longrightarrow -4\lambda h.
\end{equation}
Consequently \eqref{forward} tends to
\begin{equation}\label{dycklimit}
 F(w)=\left(\prod_{\text{annihilations}}h\right)
 q^{\sum_{\text{creations}}h}
 =\prod_{\text{annihilations}}h q^{h-1}.
\end{equation}
The two height sums agree because each upward edge of a Dyck path
has a downward edge at the same level. The shift \eqref{shift} has
exactly these moments. Together with \eqref{backbound}, this proves
Theorem~\ref{wordthm} for all fixed words, not only creation norms.

The creation norms themselves are exact at finite size:
\begin{equation}\label{creationnorm}
 s_{k,N}=\phi_{O,N}\bigl(M_N^k(M_N^\dagger)^k\bigr)
 =k!\prod_{j=0}^{k-1}R_N(j)
 \longrightarrow k!q^{k(k-1)/2}=:s_k.
\end{equation}
There is no backward matching in this ordering. In particular,
$\phi(M^2(M^\dagger)^2)=2q$, whereas
$\phi(MM^\dagger MM^\dagger)=1$. These distinguish the dressed mode
from a bare canonical oscillator.

\newpage
\section{Canonical variable and the new vacuum representation}
On the standard bosonic finite-particle core, with
$b|n\rangle=\sqrt n|n-1\rangle$ and $\nn=b^\dagger b$, the shift is
\begin{equation}\label{dress}
 M=q^{\nn/2}b,\qquad M^\dagger=b^\dagger q^{\nn/2}.
\end{equation}
The order matters: $M$ acts with $q^{(n-1)/2}$, not $q^{n/2}$.
Conversely, on that core,
\begin{equation}\label{undress}
 b=q^{-\nn/2}M,\qquad b^\dagger=M^\dagger q^{-\nn/2},
 \qquad [b,b^\dagger]=I.
\end{equation}
The inverse dressing is unbounded. Equation \eqref{undress} is first
a statement about the limiting representation and its dense domain.
There is also an actual microscopic realization on fixed excitation
levels, as follows.

The bare canonical/dressed distinction already appears in
\cite[HTML Eqs.~(75), (85)--(86)]{RX}; the corresponding PDF equations
are (3.43), (3.53)--(3.54). Here the single-mode law was obtained from
\eqref{ham}--\eqref{block}, not from the Majorana mixed-chord rules.

\subsection*{An actual finite-size canonical sequence}
Let $Q_E=\sum_{i\in E}c_i^\dagger c_i$, $n_N=Q_E/p$, and
$q_N=e^{-4\lambda_N}$. Define finite-size operators
\begin{equation}\label{microb}
 b_N=q_N^{-n_N/2}M_N,\qquad
 b_N^\dagger=M_N^\dagger q_N^{-n_N/2}.
\end{equation}
They are adjoints, and $[n_N,M_N^\dagger]=M_N^\dagger$ exactly.
Every intermediate state obtained from $\Om$ along a word has
$n_N$ equal to that word's height. This remains true when $M_N$
has components outside the span of pure-creation vectors.
For a Dyck word define $A(w)=\sum_{\rm annihilation}(h-1)$.
Each creation contributes $q_N^{-h/2}$ and each annihilation
$q_N^{-(h-1)/2}$, so, sample by sample,
\begin{equation}\label{microword}
 \langle O|w(b_N,b_N^\dagger)|O\rangle
 =q_N^{-A(w)}\langle O|w(M_N,M_N^\dagger)|O\rangle.
\end{equation}
Theorem~\ref{wordthm} then yields all fixed canonical vacuum moments:
\begin{equation}
 \lim_N\phi_{O,N}(w(b_N,b_N^\dagger))
 =\prod_{\rm annihilation}h.
\end{equation}
Non-Dyck or unbalanced words remain zero. Expanding a squared norm
in finitely many such words proves, for every fixed noncommutative
polynomial $P$,
\begin{equation}\label{coreccr}
 \EE_J\left\|([b_N,b_N^\dagger]-I)
        P(b_N,b_N^\dagger)\Om\right\|^2\longrightarrow0.
\end{equation}
Thus the canonical sequence is microscopic, with convergence on the
specified Slater polynomial core. The norm of the inverse dressing
on the full Fock space grows with $N$; no uniform operator-norm
limit or exact finite-matrix CCR is asserted.

\newpage
\subsection*{Actual ladder vectors and the limiting algebra}
For fixed $k$ and sufficiently large $N$ such that $kp\le\min(a,m)$, normalize
$|k;N\rangle=(M_N^\dagger)^k\Om/\sqrt{s_{k,N}}$ in the annealed
state/disorder Hilbert space. Different $k$ are orthogonal by their
particle--hole excitation numbers. Creation maps one to the next
exactly, with coefficient $\sqrt{s_{k+1,N}/s_{k,N}}$.
The all-word theorem gives
\begin{equation}\label{strongladder}
 \left\|M_N|k;N\rangle
       -\sqrt{s_{k,N}/s_{k-1,N}}\,|k-1;N\rangle\right\|^2
 \longrightarrow0\qquad(k\ge1).
\end{equation}
Indeed, the overlap with $|k-1;N\rangle$ equals
$\sqrt{s_{k,N}/s_{k-1,N}}$ exactly, and the norm squared of the
first term tends to $kq^{k-1}$ by a fixed word of length $2k+2$.
Their difference therefore vanishes. This is fixed-level annealed
norm convergence, not operator-norm convergence on the entire
microscopic Hilbert space.

\paragraph{Algebra generated by the adapted mode.}
The weights in \eqref{shift} are positive and tend to zero; $M$ is a
compact bounded operator. Its kernel is exactly $\mathbb C|0\rangle$,
so the vacuum projection $P_0=1_{\{0\}}(M^\dagger M)$ belongs to
$W^*(M)$. The operators $(M^\dagger)^nP_0M^j$ are nonzero multiples
of every matrix unit $|n\rangle\langle j|$. Thus
\begin{equation}
 W^*(M)=B(\ell^2(\mathbb N_0)).
\end{equation}
The vacuum is cyclic and not separating. This is a new adapted
algebra; it does not split the old tracial field algebra into physical
annihilators. The separating tracial chord vacuum of \cite{Xu} refers
to its specified field algebra, not all auxiliary Fock ladders.

The difference of states is already visible at finite size. Averaging
the Gaussian coefficients and taking the normalized full-Fock trace
gives
\begin{equation}\label{invisible}
 \EE_J\tr_{\rm full}(M_N^\dagger M_N)=4^{-p},\qquad
 \phi_{O,N}(M_NM_N^\dagger)=1.
\end{equation}
Each diagonal term imposes $p$ occupied and $p$ empty sites, so its
normalized trace is $2^{-2p}$. The same microscopic sequence therefore
vanishes in the old tracial $L^2$ norm and survives in the adapted
state. It is not an observable identified by a state-preserving map
from the old tracial limiting algebra.

\paragraph{Fixed scalar words also concentrate.}
This is not needed for the annealed theorem, but follows from the same
label argument. For $F_{w,N}(J)=\langle O|w(M_N,M_N^\dagger)|O\rangle$
of length $2k$, expand the variance in two vacuum copies. Wick
matchings internal to each copy cancel. In every remaining matching
a shared label appears only once in each copy. All of its sites
must occur in the other labels, or that copy cannot return to the
vacuum. Conditioning on at most $2k-1$ other labels gives
\begin{equation}
 \operatorname{Var}F_{w,N}\le (2k)!\rho_N(2k-1)\longrightarrow0.
\end{equation}
For the microscopic $b_N$ words, \eqref{microword} multiplies this
bound by a fixed-word factor tending to $q^{-2A(w)}$.
No concentration claim for the joint Hamiltonian dynamics is needed
or inferred here.

\newpage
\section{The full Hamiltonian: a Poisson family of missing directions}
The adapted mode captures an actual part of $H_N\Om$. To measure the
rest, let $k$ be the number of particles placed in $E$ by a term of
$H_N$ (and thus the number of holes in $O$). The exact annealed squared
norm in this sector is
\begin{equation}\label{weights}
 W_{k,N}=w_N\binom mp\binom pk\binom ak,
 \qquad 0\le k\le p.
\end{equation}
To see this, $B$ must be an occupied $p$-set. After removing it,
the creator set $A$ can fill $p-k$ of its $p$ holes and $k$ of the
$a$ empty sites. Squaring and averaging pairs a term only with its
Hermitian reverse; no independent-probe approximation is made.
Vandermonde's identity gives
\begin{equation}\label{cm}
 C_{m,N}=\phi_{O,N}(H_N^2)
 =w_N\binom mp\binom{a+p}p\longrightarrow C=J_s^2/\lambda.
\end{equation}

For the deficit $r=p-k$, the normalized weights and their falling
factorial moments are
\begin{align}
 \Pr_N(r)&=\frac{\binom pr\binom a{p-r}}{\binom{a+p}p},\label{deficit}\\
 \EE_N\ff rt&=\frac{\ff pt^2}{\ff{a+p}t}
       \longrightarrow(2\lambda)^t.\label{factmom}
\end{align}
This is a hypergeometric law: choose $p$ items from $a+p$, of which
$p$ are marked. For $t\le p$, its factorial moments are bounded by
$(p^2/a)^t$. Hence for every fixed $z\ge1$,
$\EE_N z^r\le\exp[(z-1)p^2/a]$.
The factorial-moment expansion is uniformly dominated, identifying
the limit as Poisson$(2\lambda)$. In particular,
\begin{equation}\label{captured}
 \frac{\sigma_N^2}{C_{m,N}}
 =\frac{\binom ap}{\binom{a+p}p}
 \longrightarrow e^{-2\lambda}=\sqrt q.
\end{equation}
This is a squared-norm fraction. The normalized overlap amplitude
between $H_N\Om$ and $M_N^\dagger\Om$ tends to
$e^{-\lambda}=q^{1/4}$.

Let $\KK_N$ be the closed span of $(M_N^\dagger)^j\Om$ in the
annealed joint Hilbert space. At finite $N$, $H_N\Om$ overlaps only
its one-block vector: the mean vacuum overlap vanishes and higher
vectors have more than $p$ excited particles. Therefore exactly
\begin{align}
 P_{\KK_N}H_N\Om&=\sigma_N M_N^\dagger\Om,\label{project}\\
 \|(1-P_{\KK_N})H_N\Om\|^2
 &=C_{m,N}-\sigma_N^2\longrightarrow C(1-\sqrt q)>0.\label{leak}
\end{align}
The physical $H_N$ does not close on this single ladder. A guessed
number Hamiltonian would lose the nonzero norm in \eqref{leak}.

\newpage
\section{Exact checks and the next constructive calculation}
The new driver \texttt{code/check\_slater\_ladder.py} uses integer CAR
actions and rational Gaussian Wick sums. Its first successful output
is preserved in the stage's \texttt{exact-checks.json}; it was not
rerun to overwrite an earlier result.
\begin{center}
\begin{tabular}{lr}
\toprule
Finite fixtures $(N,p;k_{\max})$ & $(4,1;4),(6,1;3),(6,2;3),(8,2;3)$\\
\midrule
Gaussian label assignments & 1,537,310\\
Forward CAR/active-interval comparisons & 754,030\\
Complete Dyck-word moments & 46\\
Non-Dyck words vanishing by right action & 136\\
Words with nonzero backward corrections & 25\\
Exact hypergeometric diagnostic rows & 6\\
Mixed-$Q$ linear-cumulant checks & 189\\
\bottomrule
\end{tabular}
\end{center}
The fixtures also enumerate every original $p$-body monomial in
$H_N\Om$, independently checking \eqref{weights} and \eqref{cm}.
The hypergeometric rows use $\lambda_N=1/2$, $p=2,4,8,16,32,64$,
$N=2p^2$. Saved total-variation distances to Poisson$(1)$ are
floating-point diagnostics; the exact factorial moments and the
analytic bound above supply the convergence argument. Small fixtures
test identities and signs; they do not prove the double-scaling limit.

The next microscopic object is now specified: preserve the same
$J$ and compute the joint ordered moments of $H_N,M_N,M_N^\dagger$,
including the $r\ge1$ sectors of \eqref{deficit}. The first target is a
positive representation that reproduces \eqref{project}--\eqref{leak}
and the interleaved fourth moments. A subsequent step must connect
this adapted algebra to the already constructed singlet bilocal
observables. The chosen Slater state is not stationary, so its
preparation and real-time evolution are part of that calculation.

Only after these steps can one test a space/boost generator and a
meson mass operator against the proposed bulk theory. The reference
polarization needed for canonical bilocals in light-front QCD is
explicit in Itakura and Rajeev \cite{Itakura,Rajeev}; it supplies a
useful comparison, not an identification of our block number with
their momentum or mass variables. Neither a finite-$G$ interaction
nor a complete de Sitter dictionary has been derived here.

\paragraph{Reproducibility.}
The JSON records parameters, versions, parameter digest, driver hash
and parent-manifest hashes. The local evidence recorder additionally
binds sources, proofs, manuscript, two compiler passes and all final
page renders. It verifies preserved files and saved arithmetic; it
does not certify a theorem. The prior bilocal scientific outputs,
including historical failures in the separate KS branch, are retained.

\newpage
\appendix
\section{A supporting check on the old tracial low fields}
The preceding three-species law has normalized self-adjoint fields
$(h,u,x)$ with covariance $I$ and
\begin{equation}
 Q=\begin{pmatrix}q&q&q^2\\q&q&q^2\\q^2&q^2&q^4\end{pmatrix}.
\end{equation}
For real $Y=\alpha h+\beta u+\gamma x$, set
$A=\alpha^2+\beta^2$, $B=\gamma^2$. Its fourth classical cumulant is
the direct mixed-$Q$ Wick consequence
\begin{equation}\label{cum}
 \tau(Y^4)-3\tau(Y^2)^2
 =-(1-q)A^2-2(1-q^2)AB-(1-q^4)B^2.
\end{equation}
It is strictly negative for every nonzero $Y$ at fixed $0<q<1$,
and lies between $(q^4-1)(A+B)^2$ and $(q-1)(A+B)^2$.
Thus no nonzero linear whitening of this triple is a Gaussian field
in a quasifree canonical state. This is an elementary implication of
the previously obtained law, not a new mixed-$Q$ theorem and not an
obstruction to every nonlinear or nonquasifree representation.
It motivates retaining explicit polarization rather than only
changing a low-field Gram matrix.

The companion KMS note separately treats joint annealed Gibbs
tilting, exact thermofield realization, and smooth frequency filters.
Its Gaussian dilation matches two-point functions only. None of
these constructions is substituted for the same-$J$ Slater derivation
in the main text.

\begin{thebibliography}{9}
\bibitem{MSS} S. Miyashita, Y. Sekino and L. Susskind,
\emph{Holograms and Standard Models},
\href{https://arxiv.org/abs/2607.05678v2}{arXiv:2607.05678v2} (2026).
\bibitem{RX} H. Rajgadia and J. Xu,
\emph{Emergent States and Algebras from the Double-Scaling limit of
Pure States in SYK},
\href{https://arxiv.org/abs/2604.14387v3}{arXiv:2604.14387v3} (2026).
\bibitem{Xu} J. Xu, \emph{Von Neumann Algebras in Double-Scaled SYK},
\href{https://arxiv.org/abs/2403.09021v7}{arXiv:2403.09021v7} (2026 revision).
\bibitem{Itakura} K. Itakura,
\emph{Boson Expansion Methods in (1+1)-dimensional Light-Front QCD},
\href{https://arxiv.org/abs/hep-th/9604032v2}{arXiv:hep-th/9604032v2} (1996).
\bibitem{Rajeev} S. G. Rajeev, \emph{Quantum Hadrondynamics in Two Dimensions},
\href{https://arxiv.org/abs/hep-th/9401115}{arXiv:hep-th/9401115} (1994).
\end{thebibliography}
\end{document}
